Die Gliederung der Arbeit erfolgt mit den gewöhnlichen
\LaTeX-Gliederungsbefehlen \verb+\chapter+, \verb+\section+,
\verb+\subsection+, usw. Vergleiche hierzu auch die
Listings~\ref{lst:tpl:article} bis~\ref{lst:tpl:book}. Die \texttt{*}-Varianten
werden nicht nummeriert und erscheinen standardmäßig nicht im
Inhaltsverzeichnis. Soll ein unnummierierten Punkt dennoch ins
Inhaltsverzeichnis aufgenommen werden, kann dies über den
\verb+\addcontentsline+-Befehl erfolgen, wie in
Listing~\ref{lst:tpl:addcontentsline} gezeigt.\todo{Bib\TeX}\ 
Der Befehl muss immer nach dem entsprechenden Gliederungsbefehl aufgerufen
werden. Sonst stimmen die Seitenverweise nicht.

\lstinputlisting[language={[LaTeX]{TeX}},style=block,escapechar=\!,float={htb},caption={Aufnahme zusätzlicher Verzeichniseinträge},label={lst:tpl:addcontentsline},morekeywords={chapter,subsection}]{code/tpl_addcontentsline}

Für die Gliederung ist folgendes Musterschema zweckmäßig. Je nach Art der
Arbeit entfallen die nicht zutreffenden Punkte.

\begin{compactitem}
\item Deckblatt
\item Widmung
\item Kurzreferat
\item Aufgabenstellung
\item Inhaltsverzeichnis
\item weitere Verzeichnisse
      (ggf. auf am Ende der Arbeit möglich, siehe unten)
\item Abkürzungsverzeichnis, Symbole und Bezeichner
\item Textkörper
\item Anhänge
\item Literaturverzeichnis
\item weitere Verzeichnisse
      (Alternative, sofern nicht bereits am Anfang gesetzt)
\item Glossar
\item Index
\item Danksagung
\item Selbständigkeitserklärung
\end{compactitem}
